\section{Escape below the critical scale}\label{sec:escape}
The localized disk estimate must exclude a favorite at every time in a block,
not just at its endpoints. Monotonicity of local times makes this possible:
use the global maximum at the left endpoint and the local time in the target
disk at the right endpoint.

\begin{proposition}[Escape below the critical scale]\label{prop:supercritical-escape}
For every \(\gamma>\tfrac12\), almost surely,
\begin{equation}
\liminf_{n\to\infty}
       \frac{R_n^-}{\sqrt n/(\log n)^\gamma}=\infty.          \tag{80}
\end{equation}
\end{proposition}

\begin{proof}
Fix \(\gamma>1/2\) and a positive integer spatial multiplier \(A_0\).
Choose
\[
       0<\delta<\frac{2\gamma-1}{4}.
\]
Set \(N_j=\lfloor e^{2j}\rfloor\).
Proposition~\ref{prop:deterministic-maximum-lower} gives, for all
sufficiently large \(j\),
\begin{equation}
M_{N_j}\ge u_{j,2+\delta}.                             \tag{81}
\end{equation}

Consider any \(n\in[N_j,N_{j+1}]\cap\mathbb N\).  If
\(x\in\mathcal F_n\) lies in
\[
       B\left(0,A_0\frac{\sqrt n}{(\log n)^\gamma}\right),
\]
then monotonicity of local times and (81) give
\[
 L_{N_{j+1}}(x)\ge L_n(x)=M_n\ge M_{N_j}
       \ge u_{j,2+\delta}.
\]
Since \(n\le N_{j+1}\) and \(\log n\ge\log N_j\), the same site lies in
\begin{equation}
B\left(0,C_{A_0}e^{j+1}j^{-\gamma}\right)             \tag{82}
\end{equation}

Take
\[
       K_j=C_*e^{j+1}\sqrt{\log j},\qquad J'_j=\log K_j.
\]
By the maximal inequality in Lemma~\ref{lem:elementary-walk-estimates}, \(C_*\)
can be chosen so that
\begin{equation}
\sum_jP(\tau_{K_j}\le N_{j+1})<\infty.                 \tag{83}
\end{equation}
For large \(j\), the ball (82) is contained in
\[
       B(0,K_j(J'_j)^{-\gamma}).
\]
Moreover \(J'_j=j+O(\log\log j)\), and the fixed extra coefficient
\(\delta\) gives
\begin{equation}
u_{j,2+\delta}\ge u_{J'_j,2+2\delta}                  \tag{84}
\end{equation}
for all large \(j\): the negative term
\(-\delta g j\log j\) dominates the \(O(j\log\log j)\) change in the
leading term.

Choose independently an upper-error exponent \(A_{\rm err}>2\).
On the complement of the summable event in (83), the implication
(81)--(84) and Proposition~\ref{prop:localized-disk-upper}, with
\(\beta=\gamma\) and \(c=2+2\delta\), show that the probability of any
such \(n\) in the whole block is at most
\begin{equation}
(J'_j)^{-2\gamma+2\delta+o(1)}+O(j^{-A_{\rm err}}).     \tag{85}
\end{equation}
Our choice of \(\delta\) makes the power in (85) strictly smaller than
\(-1\).  Hence (85) is summable.  Borel--Cantelli proves that eventually
no favorite in the \(j\)-th block lies in the ball with multiplier \(A_0\).
Intersecting over \(A_0\in\mathbb N\) proves (80).
\end{proof}


\begin{proof}[Proof of Theorem~\ref{main:precise-escape-rate}]
Propositions~\ref{prop:critical-recurrence} and
\ref{prop:supercritical-escape} give the two assertions for each fixed
exponent. Intersect the probability-one events for rational
$\gamma>1/2$ and for the single critical exponent $1/2$.
If $\gamma>1/2$ is real, choose a rational $q\in(1/2,\gamma)$; then
\[
 \frac{R_n^-}{a_\gamma(n)}
 =\frac{R_n^-}{a_q(n)}(\log n)^{\gamma-q}\longrightarrow\infty.
\]
For $\gamma\le1/2$, the inequality
$R_n^+/a_\gamma(n)\le R_n^+/a_{1/2}(n)$ gives the other assertion on the
same event. This proves the simultaneous formulation.
\end{proof}
